% Part: computability
% Chapter: computability-theory
% Section: complement-ce

\documentclass[../../../include/open-logic-section]{subfiles}

\begin{document}

\olfileid{cmp}{thy}{cmp}
\olsection{સંગણકીય રીતે પરિગણનીય ગણો પૂરક હેઠળ બંધ નથી}

માનો કે $A$ સંગણકીય રીતે પરિગણનીય છે. શું $A$નો પૂરક,
$\Complement{A} = \Nat \setminus A$, હંમેશાં સંગણકીય રીતે
પરિગણનીય હોય છે? આગળનું પ્રમેય અને પરિણામ બતાવે છે કે
જવાબ ``ના'' છે.

\begin{thm}
\ollabel{thm:ce-comp}
$A$ પ્રાકૃતિક સંખ્યાઓનો કોઈ પણ ગણ હોય. ત્યારે $A$ સંગણનીય
ત્યારે અને માત્ર ત્યારે છે જ્યારે $A$ અને $\Complement{A}$
બંને સંગણકીય રીતે પરિગણનીય હોય.
\end{thm}

\begin{proof}
એક દિશા સહેલી છે: $A$ સંગણનીય હોય તો $\Complement{A}$ પણ
સંગણનીય છે ($\Char{A} = 1 \tsub
\Char{\Complement{A}}$), તેથી બંને સંગણકીય રીતે પરિગણનીય છે.

ઊલટી દિશામાં, માનો કે $A$ અને~$\Complement{A}$ બંને
સંગણકીય રીતે પરિગણનીય છે. $A$, $\cfind{d}$નો પ્રદેશ
અને $\Complement{A}$, $\cfind{e}$નો પ્રદેશ હોય. $h$ને
આ રીતે વ્યાખ્યાયિત કરો:
\[
h(x) = \umin{s}{(T(d,x,s) \lor T(e,x,s))}.
\]
બીજા શબ્દોમાં, ઇનપુટ~$x$ પર $h$, $\cfind{d}$ અથવા
$\cfind{e}$ની થંભતી સંગણના શોધે છે. હવે $x \in A$ હોય
તો પહેલી સ્થિતિમાં અને $x \in \Complement{A}$ હોય તો
બીજી સ્થિતિમાં શોધ સફળ થાય છે. તેથી $h$ સંપૂર્ણ સંગણનીય
વિધેય છે. પરંતુ હવે દરેક~$x$ માટે $x \in A$ ત્યારે અને
માત્ર ત્યારે જ્યારે $T(d, x, h(x))$, એટલે વ્યાખ્યાયિત
થનાર વિધેય $\cfind{d}$ હોય. $T(d, x, h(x))$ સંગણનીય
સંબંધ હોવાથી $A$ સંગણનીય છે.
\sourcecorrection{OLCT-006}{સ્થિર અંગ્રેજી મૂળ આ અંતિમ
પરિક્ષણમાં $T(e, x, h(x))$ અને $\cfind{e}$ લખે છે,
જ્યારે ઉપર $d$નો પ્રદેશ $A$ અને $e$નો પ્રદેશ તેનો પૂરક
રાખ્યો છે. તેથી $A$ માટે અહીં બંને જગ્યાએ $d$ જોઈએ.}
\end{proof}

\begin{explain}
અનૌપચારિક સંગણનાની ભાષામાં વાત વધુ સહેલી છે: $A$નો નિર્ણય
કરવા ઇનપુટ $x$ પર $\cfind{d}$ અને $\cfind{e}$ની થંભતી
સંગણનાઓ શોધો. બેમાંથી એક તો જરૂર થંભશે. જો તે
$\cfind{d}$ હોય, તો $x$ એ~$A$માં છે; અન્યથા
$x$, $\Complement{A}$માં છે.
\sourcecorrection{OLCT-007}{સ્થિર અંગ્રેજી મૂળ આ
સમજૂતીમાં $\cfind{e}$ અને $\cfind{f}$ લખીને પહેલા
વિધેયને $A$ માટે ગણે છે. ઉપર નિર્ધારિત સૂચકો $d$ અને
$e$ છે; $A$નો પ્રદેશ $\cfind{d}$નો છે.}
\end{explain}

\begin{cor}
\ollabel{cor:comp-k}
$\Complement{K_0}$ સંગણકીય રીતે પરિગણનીય નથી.
\end{cor}

\begin{proof}
આપણે જાણીએ છીએ કે $K_0$ સંગણકીય રીતે પરિગણનીય છે,
પરંતુ સંગણનીય નથી. જો $\Complement{K_0}$ સંગણકીય રીતે
પરિગણનીય હોત, તો \olref{thm:ce-comp} મુજબ $K_0$
સંગણનીય હોત, જે \olref[ncp]{thm:K-0} સાથે વિરોધાભાસ છે.
\end{proof}

\end{document}
